El sitio oficial del Concurso Internacional Universitario de Programación (ICPC, por sus siglas en inglés) menciona que la ICPC es ``un concurso de programación algoritmica para estudiantes universitarios'', aunque esta definición es correcta, si el lector nunca ha participado en esta, es probable que no tenga clara cuál es la diferencia entre esta competencia y otras que existen en el mundo de la programación, por ejemplo hackatones famosos como ``Microsoft Global Hackaton'' o ``NASA Space Apps'', en los cuáles se pone enfásis en desarrollar soluciones a problemas también. La diferencia radica principalmente en que en las últimas se pone mucho enfásis en desarrollar proyectos completos o aplicaciones durante varios días que resuelven un gran problema, generalmente del mundo real y abierto a que los concursantes exploren diferentes perspectivas del mismo, por ejemplo en ``NASA Space Apps 2025'' uno de los problemas a resolver era el de buscar exoplanetas usando inteligencia artificial.

Por otro lado en la ICPC y otras olimpiadas de programación lo que se busca es que los concursantes programen soluciones a problemas concretos, en los que la correctez de una solución se basa en si esta arroja la respuesta correcta a todos los casos de prueba propuestos, esto es evaluado por un juez, el cuál se explicará cómo funciona en el siguiente capítulo.

Las competencia más importante de programación competitiva a nivel universitario actualmente es la ICPC, esta es una competencia por equipos de 3 personas \footnote{Estos deben ser estudiantes universitarios de la misma institución}, que se realiza de forma anual con varias fases, en caso de México, desde el año 2024 se introdujo una nueva fase, de forma que actualmente consta de 4 etapas. La primera etapa es el Gran Premio de México, el cuál consiste en 3 fechas, en las que en cada una se dejan alrededor de 9 a 13 problemas que como se verá en este texto abarcan diferentes temas. Los equipos tienen 5 horas para poder resolver la mayor cantidad de problemas posible, es importante mencionar que si 2 o más equipos resuelven la misma cantidad de problemas el orden de la tabla se decidirá por qué equipo tiene menor penalización, la cuál es la suma de los tiempos, en minutos, en los cuales se resolvió cada problema, cada intento incorrecto sumará 20 minutos extra de penalización al equipo en caso de resolver el problema del que fue el envío incorrecto. Una vez que concluye cada fecha, los puntos se otorgan como se muestra en la tabla de abajo a los primeros cincuenta lugares y a partir del 51, todos los equipos que resuelvan mínimo un problema reciben 1 punto.
\begin{adjustwidth}{-2.5 cm}{-2.5 cm}\centering\begin{threeparttable}
\caption{Generated by Spread-LaTeX}\label{tab: }
\begin{tabular}{lrrrr}\toprule
\textbf{Lugares} &\textbf{Primera Fecha} &\textbf{Segunda Fecha} &\textbf{Tercera Fecha} \\\midrule
Lugar 1 &100 puntos &150 puntos &200 puntos \\
Lugar 2 &95 puntos &145 puntos &195 puntos \\
Lugar 3 &90 puntos &140 puntos &190 puntos \\
Lugar 4 &88 puntos &135 puntos &185 puntos \\
Lugar 5 &86 puntos &130 puntos &180 puntos \\
Lugar 6 &84 puntos &125 puntos &175 puntos \\
Lugar 7&82 puntos &120 puntos &170 puntos \\
Lugar 8&80 puntos &115 puntos &165 puntos \\
Lugar 9&78 puntos &110 puntos &160 puntos \\
Lugar 10&76 puntos &105 puntos &155 puntos \\
Lugar 11 &74 puntos &100 puntos &150 puntos \\
Lugar12 &72 puntos &95 puntos &145 puntos \\
Lugar 13 &70 puntos &90 puntos &140 puntos \\
Lugar 14 &68 puntos &85 puntos &135 puntos \\
Lugar 15 &66 puntos &80 puntos &130 puntos \\
Lugar 16 &64 puntos &75 puntos &125 puntos \\
Lugar 17 &62 puntos &70 puntos &120 puntos \\
Lugar 18 &60 puntos &65 puntos &115 puntos \\
Lugar 19 &58 puntos &63 puntos &110 puntos \\
Lugar 20 &56 puntos &61 puntos &105 puntos \\
Lugar 21 &54 puntos &59 puntos &100 puntos \\
Lugar 22 &52 puntos &57 puntos &95 puntos \\



\bottomrule
\end{tabular}
\end{threeparttable}\end{adjustwidth}
\begin{adjustwidth}{-2.5 cm}{-2.5 cm}\centering\begin{threeparttable}
\begin{tabular}{lrrrr}\toprule
\textbf{Lugares} &\textbf{Primera Fecha} &\textbf{Segunda Fecha} &\textbf{Tercera Fecha} \\\midrule
Lugar 23 &50 puntos &55 puntos &90 puntos \\
Lugar 24 &48 puntos &53 puntos &85 puntos \\
Lugar 25 &46 puntos &51 puntos &80 puntos \\
Lugar 26 &44 puntos &49 puntos &75 puntos \\
Lugar 27 &42 puntos &47 puntos &70 puntos \\
Lugar 28 &40 puntos &45 puntos &65 puntos \\
Lugar 29 &38 puntos &43 puntos &60 puntos \\
Lugar 30 &36 puntos &41 puntos &55 puntos \\
Lugar 31 &34 puntos &39 puntos &50 puntos \\
Lugar 32 &32 puntos &37 puntos &45 puntos \\
Lugar 33 &30 puntos &35 puntos &40 puntos \\
Lugar 34 &28 puntos &33 puntos &35 puntos \\
Lugar 35 &26 puntos &31 puntos &33 puntos \\
Lugar 36 &24 puntos &29 puntos &31 puntos \\
Lugar 37 &22 puntos &27 puntos &29 puntos \\
Lugar 38 &20 puntos &25 puntos &27 puntos \\
Lugar 39 &18 puntos &23 puntos &25 puntos \\
Lugar 40 &16 puntos &21 puntos &23 puntos \\
Lugar 41 &14 puntos &19 puntos &21 puntos \\
Lugar 42 &12 puntos &17 puntos &19 puntos \\
Lugar 43 &10 puntos &15 puntos &17 puntos \\
Lugar 44 &8 puntos &13 puntos &15 puntos \\
Lugar 45 &7 puntos &11 puntos &13 puntos \\
Lugar 46 &6 puntos &9 puntos &11 puntos \\
Lugar 47 &5 puntos &7 puntos &9 puntos \\
Lugar 48 &4 puntos &5 puntos &7 puntos \\
Lugar 49 &3 puntos &3 puntos &5 puntos \\
Lugar 50 &2 puntos &2 puntos &3 puntos \\

\bottomrule
\end{tabular}
\end{threeparttable}\end{adjustwidth}

Tras las 3 fechas los cupos para clasificar al regional, también llamado ``ICPC Mexico Finals'' se asignan por varias categorías, es importante mencionar que la cantidad de cupos de cada tipo así como sus reglas suelen variar año con año, se recomienda consultar el sitio oficial constantemente, las categorías actuales de cupos para el regional de 2025 son:

\begin{itemize}
    \item Cupos por desempeño: Estos son los 44 equipos mejor posicionados en la tabla al terminar las primeras dos fechas.
    \item Cupo por institución sede: La institución que sea sede del regional tendrá un cupo extra para un equipo de la institución que no haya clasificado por desemepño.
    \item Cupo rosa: Es un cupo que se le da al mejor no clasificado por los criterios anteriores en el que todas las integrantes sean mujeres.
    \item Cupo rosa sede: Este cupo es dado a la institución sede del regional para el mejor equipo de 3 mujeres de la institución que no haya clasificado.
    \item Cupos comunidades: Son dos cupos a las comunidades con mayor participación en la competencia.
    \item Cupos repechaje: La tercera de las fechas se le llama repechaje y es una oportunidad extra para los equipos de clasificar a la ``ICPC Mexico Finals'', se otorgan 6 lugares a los mejores equipos en la tercera fecha sin contar equipos clasificados por otros criterios.
\end{itemize}






%terminar de explica:
%-cuantos puntos dan a cada equipo
%-al final de las 3 fechas la cantidad de cupos (rpechaje, desempeño, rosa etc)
%regional
%latam championship
%mundial
%Agregar ahora sí problemas de ejemplo
